\subsection{Стандартный линейный вязко-упругий материал}

Две модели уже построены:
\[
  \text{(1) Кельвин–Фойгт:}\quad \Sdev = 2\mu\,\dev + 2\eta\,\dt{\dev} ,
  \qquad
  \text{(2) Максвелл:}\quad \dt{\dev} = \frac{\dt{\Sdev}}{2\mu} + \frac{\Sdev}{2\eta} .
\]
Хотим их объединить — так, чтобы модель имела и ползучесть, и релаксацию.

\inputfig{figures/fig-standard-combine}

\subsubsection*{Вывод определяющего соотношения}

Ветвь Максвелла ($\mu$, $\eta$) параллельно пружине $\mu_1$. По аксиомам
параллельного соединения
\[
  \Sdev = \Sdev_1 + \Sdev_2 ,
  \qquad
  \dev = \dev_1 = \dev_2 ,
\]
\[
  \Sdev_1 = 2\mu_1 \dev ,
  \qquad
  \dt{\dev}_2 = \frac{\dt{\Sdev}_2}{2\mu} + \frac{\Sdev_2}{2\eta} ,
  \qquad
  \Sdev_2 = \Sdev - \Sdev_1 = \Sdev - 2\mu_1 \dev .
\]
Подставляя $\dt{\dev} = \dt{\dev}_2$:
\[
  \dt{\dev}
  = \frac{\dt{\Sdev} - 2\mu_1 \dt{\dev}}{2\mu}
  + \frac{\Sdev - 2\mu_1 \dev}{2\eta}
\]
и, собирая слева всё с $\dev$:
\[
  \left(1 + \frac{\mu_1}{\mu}\right)\dt{\dev}
  + \frac{\mu_1}{\eta}\,\dev
  = \frac{\dt{\Sdev}}{2\mu} + \frac{\Sdev}{2\eta} .
\]
Поделив на $2\mu_1$, получаем симметричную форму, записанную на доске:
\[
  \boxed{\ \left(\frac{1}{2\mu_1} + \frac{1}{2\mu}\right)\dt{\dev}
         + \frac{1}{2\eta}\,\dev
         = \frac{1}{2\mu_1}\left(\frac{\dt{\Sdev}}{2\mu} + \frac{\Sdev}{2\eta}\right)\ }
\]

\subsubsection*{Что задаём — то и получаем}

В зависимости от того, что задано, из одного и того же уравнения получается
ползучесть или релаксация:
\[
  \begin{aligned}
    \text{задаём } \Sdev &\ \Longrightarrow\ \text{получаем (1) — ползучесть} , \\
    \text{задаём } \dev &\ \Longrightarrow\ \text{получаем (2) — релаксацию} .
  \end{aligned}
\]

\subsubsection*{Предельные случаи}

\[
  \mu_1 \to 0:\quad
  \dt{\dev} = \frac{\dt{\Sdev}}{2\mu} + \frac{\Sdev}{2\eta}
  \quad \text{— Максвелл;}
\]
\[
  \mu \to \infty:\quad
  \dt{\dev} + \frac{\mu_1}{\eta}\dev = \frac{\Sdev}{2\eta}
  \ \Longleftrightarrow\
  \Sdev = 2\mu_1 \dev + 2\eta\,\dt{\dev}
  \quad \text{— Кельвин–Фойгт.}
\]

\begin{idea}
Отличие стандартного тела от исходных моделей: в отличие от Кельвина–Фойгта оно
даёт мгновенный упругий отклик на скачок нагрузки, а в отличие от Максвелла —
конечную, а не неограниченно растущую деформацию при постоянном напряжении.
Поэтому в нём есть и ползучесть, и релаксация одновременно.
\end{idea}

\begin{warn}
На схеме у доски подписано «\dots{} ползучесть $+$ \dots{} релаксация» — часть
подписи неразборчива. Вторая схема (пружина последовательно со звеном
Кельвина–Фойгта) даёт то же уравнение с другими комбинациями констант.
\end{warn}
